\chapter{Internal Categories and Mealy Morphisms in the Bicategory of Spans}
\label{ch:spans-internal-categories-mealy}

\section{Span calculus and notation}
\label{sec:span-calculus-notation}

Throughout this chapter, let \(\mathcal E\) be a category with pullbacks. We choose, once and for all, a pullback for every cospan in \(\mathcal E\). This lets us write composite spans without repeatedly specifying a choice. The bicategorical associativity and unit constraints are still the canonical isomorphisms between the corresponding chosen pullbacks.

We write composition in ordinary categorical order. Thus, if
\[
a:u\to v,
\qquad
b:v\to w,
\]
then
\[
b\circ a:u\to w
\]
means that \(a\) is performed first and \(b\) second.

A span from \(X\) to \(Y\) is written
\[
X \xleftarrow{p} S \xrightarrow{q} Y.
\]
The map \(p:S\to X\) is its left leg, and \(q:S\to Y\) is its right leg.

When no ambiguity can arise, we suppress the structure maps in pullback notation. For example, if
\[
A_1\xrightarrow{t_A}A_0
\qquad\text{and}\qquad
P\xrightarrow{p}A_0
\]
are understood, we may write
\[
A_1\times_{A_0}P
\]
for
\[
A_1\times_{t_A,A_0,p}P.
\]
When several maps into the same object are present, we display the relevant structure maps explicitly.

\section{The pullback construction of spans}
\label{sec:pullback-construction-spans}

Let
\[
X \xleftarrow{p} S \xrightarrow{q} Y
\]
and
\[
Y \xleftarrow{r} T \xrightarrow{s} Z
\]
be spans in \(\mathcal E\). Their composite is obtained by forming the pullback
\[
S\times_Y T
\]
of \(q:S\to Y\) and \(r:T\to Y\).

\[
\begin{tikzcd}
S\times_Y T \arrow[r, "\pi_T"] \arrow[d, "\pi_S"'] \arrow[dr, phantom, "\lrcorner", very near start]
&
T \arrow[d, "r"]
\\
S \arrow[r, "q"'] & Y.
\end{tikzcd}
\]

The resulting composite span from \(X\) to \(Z\) is
\[
X
\xleftarrow{p\pi_S}
S\times_Y T
\xrightarrow{s\pi_T}
Z.
\]

Equivalently:
\[
\begin{tikzcd}[column sep=large]
X
&
S\times_Y T
\arrow[l, "p\pi_S"']
\arrow[r, "s\pi_T"]
&
Z.
\end{tikzcd}
\]

A map of spans
\[
\alpha:
\left(
X \xleftarrow{p} S \xrightarrow{q} Y
\right)
\Longrightarrow
\left(
X \xleftarrow{p'} S' \xrightarrow{q'} Y
\right)
\]
is a morphism \(\alpha:S\to S'\) in \(\mathcal E\) such that
\[
p'\alpha=p,
\qquad
q'\alpha=q.
\]

This condition is represented by the commutative diagram
\[
\begin{tikzcd}[column sep=large]
& S \arrow[dl, "p"'] \arrow[dr, "q"] \arrow[d, "\alpha"] & \\
X & S' \arrow[l, "p'"] \arrow[r, "q'"'] & Y.
\end{tikzcd}
\]

\begin{definition}[Bicategory of spans]
\label{def:bicategory-of-spans}
The \textbf{bicategory of spans} in \(\mathcal E\), denoted
\[
\mathbf{Span}(\mathcal E),
\]
is the bicategory whose objects are objects of \(\mathcal E\), whose 1-cells \(X\to Y\) are spans
\[
X \xleftarrow{} S \xrightarrow{} Y,
\]
and whose 2-cells are maps of spans. The identity 1-cell on \(X\) is
\[
X \xleftarrow{1_X} X \xrightarrow{1_X} X.
\]
Horizontal composition is given by the pullback construction above.
\end{definition}

\begin{proposition}
\label{prop:span-is-bicategory}
If \(\mathcal E\) has pullbacks, then the data of Definition~\ref{def:bicategory-of-spans} form a bicategory.
\end{proposition}

\begin{proof}
Vertical composition of 2-cells is ordinary composition in \(\mathcal E\). Indeed, if
\[
\begin{tikzcd}[column sep=large]
& S \arrow[dl, "p"'] \arrow[dr, "q"] \arrow[d, "\alpha"] & \\
X & S' \arrow[l, "p'"] \arrow[r, "q'"'] \arrow[d, "\beta"] & Y \\
& S'' \arrow[ul, "p''"] \arrow[ur, "q''"'] &
\end{tikzcd}
\]
commutes, then
\[
p''\beta\alpha=p'\alpha=p,
\qquad
q''\beta\alpha=q'\alpha=q,
\]
so \(\beta\alpha:S\to S''\) is again a map of spans.

Horizontal composition of 2-cells is induced by the universal property of pullbacks. Suppose we have maps of spans
\[
\alpha:S\to S',
\qquad
\beta:T\to T',
\]
where
\[
X \xleftarrow{p} S \xrightarrow{q} Y,
\qquad
Y \xleftarrow{r} T \xrightarrow{s} Z,
\]
and
\[
X \xleftarrow{p'} S' \xrightarrow{q'} Y,
\qquad
Y \xleftarrow{r'} T' \xrightarrow{s'} Z.
\]
Since
\[
q'\alpha=q,
\qquad
r'\beta=r,
\]
we have
\[
q'\alpha\pi_S=q\pi_S=r\pi_T=r'\beta\pi_T.
\]
Hence there is a unique morphism
\[
\alpha\times_Y\beta:S\times_Y T\to S'\times_Y T'
\]
such that
\[
\pi_{S'}(\alpha\times_Y\beta)=\alpha\pi_S,
\qquad
\pi_{T'}(\alpha\times_Y\beta)=\beta\pi_T.
\]
This morphism is the horizontal composite of \(\alpha\) and \(\beta\).

Functoriality of horizontal composition, including the interchange law, follows from uniqueness in the relevant pullbacks: both sides induce the same maps on all pullback projections.

Associativity follows from the canonical isomorphism between iterated pullbacks:
\[
(S\times_Y T)\times_Z U
\cong
S\times_Y(T\times_Z U).
\]
Both sides satisfy the same universal property of triples \((s,t,u)\) satisfying the two matching equations over \(Y\) and \(Z\).

The left and right unit constraints are induced by the canonical isomorphisms
\[
X\times_X S\cong S,
\qquad
S\times_Y Y\cong S.
\]
The pentagon and triangle axioms follow from uniqueness of the canonical maps between iterated pullbacks satisfying the relevant projection equations. Therefore the data form a bicategory.
\end{proof}

\section{The endospan monoidal category}
\label{sec:endospan-monoidal-category}

Let \(C_0\) be an object of \(\mathcal E\). The hom-category
\[
\mathbf{Span}(\mathcal E)(C_0,C_0)
\]
is a monoidal category. Its objects are endospans
\[
C_0\xleftarrow{s}S\xrightarrow{t}C_0,
\]
its morphisms are maps of spans, its tensor product is horizontal composition of spans, and its tensor unit is the identity span
\[
C_0\xleftarrow{1_{C_0}}C_0\xrightarrow{1_{C_0}}C_0.
\]

We write this monoidal category as
\[
\operatorname{EndSpan}_{\mathcal E}(C_0)
:=
\mathbf{Span}(\mathcal E)(C_0,C_0).
\]

Thus, for endospans
\[
C_0\xleftarrow{s}S\xrightarrow{t}C_0,
\qquad
C_0\xleftarrow{s'}T\xrightarrow{t'}C_0,
\]
their tensor product is
\[
S\otimes T
=
S\times_{t,C_0,s'}T,
\]
regarded as the endospan
\[
C_0
\xleftarrow{s\pi_S}
S\times_{C_0}T
\xrightarrow{t'\pi_T}
C_0.
\]

\begin{remark}
\label{rem:monads-as-monoids-in-endohom}
In any bicategory \(\mathcal B\), a monad on an object \(X\) is equivalently a monoid object in the monoidal category
\[
\mathcal B(X,X),
\]
where the monoidal product is horizontal composition and the unit object is \(1_X\). Hence a monad on \(C_0\) in \(\mathbf{Span}(\mathcal E)\) is equivalently a monoid object in
\[
\operatorname{EndSpan}_{\mathcal E}(C_0).
\]
\end{remark}

\section{Endospan monoids and internal composition}
\label{sec:endospan-monoids-internal-composition}

Let
\[
C_0 \xleftarrow{s} C_1 \xrightarrow{t} C_0
\]
be an endospan. Composing this span with itself requires the pullback of
\[
t:C_1\to C_0
\qquad\text{and}\qquad
s:C_1\to C_0.
\]
We denote this pullback by
\[
C_2:=C_1\times_{t,C_0,s}C_1.
\]

\[
\begin{tikzcd}
C_2 \arrow[r, "\pi_2"] \arrow[d, "\pi_1"'] \arrow[dr, phantom, "\lrcorner", very near start]
&
C_1 \arrow[d, "s"]
\\
C_1 \arrow[r, "t"'] & C_0.
\end{tikzcd}
\]

Thus \(C_2\) is the object of pairs \((a,b)\) satisfying
\[
t(a)=s(b).
\]
The composite endospan is
\[
C_0
\xleftarrow{s\pi_1}
C_2
\xrightarrow{t\pi_2}
C_0.
\]

A monoid structure on the endospan
\[
C_0 \xleftarrow{s} C_1 \xrightarrow{t} C_0
\]
in \(\operatorname{EndSpan}_{\mathcal E}(C_0)\) consists of a unit map of spans
\[
\eta:
\left(
C_0\xleftarrow{1_{C_0}}C_0\xrightarrow{1_{C_0}}C_0
\right)
\Longrightarrow
\left(
C_0\xleftarrow{s}C_1\xrightarrow{t}C_0
\right)
\]
and a multiplication map of spans
\[
\mu:
\left(
C_0\xleftarrow{s\pi_1}C_2\xrightarrow{t\pi_2}C_0
\right)
\Longrightarrow
\left(
C_0\xleftarrow{s}C_1\xrightarrow{t}C_0
\right).
\]

Equivalently, the unit is a map
\[
e:C_0\to C_1
\]
such that
\[
se=1_{C_0},
\qquad
te=1_{C_0},
\]
and the multiplication is a map
\[
m:C_2\to C_1
\]
such that
\[
sm=s\pi_1,
\qquad
tm=t\pi_2.
\]

In generalized-element notation, if
\[
a:u\to v,
\qquad
b:v\to w,
\]
then
\[
m(a,b)=b\circ a:u\to w.
\]

\begin{theorem}[Endospan monoids are internal categories]
\label{thm:endospan-monoids-internal-categories}
Let \(\mathcal E\) be a category with pullbacks. To give a monoid object in
\[
\operatorname{EndSpan}_{\mathcal E}(C_0)
\]
is equivalently to give an internal category in \(\mathcal E\) with object-of-objects \(C_0\).
\end{theorem}

\begin{proof}
A monoid object in \(\operatorname{EndSpan}_{\mathcal E}(C_0)\) consists of an endospan
\[
C_0\xleftarrow{s}C_1\xrightarrow{t}C_0,
\]
a unit map
\[
e:C_0\to C_1
\]
satisfying
\[
se=1_{C_0},
\qquad
te=1_{C_0},
\]
and a multiplication map
\[
m:C_1\times_{t,C_0,s}C_1\to C_1
\]
satisfying
\[
sm=s\pi_1,
\qquad
tm=t\pi_2.
\]

The monoid axioms in the non-strict monoidal category
\[
\operatorname{EndSpan}_{\mathcal E}(C_0)
\]
are read after inserting the canonical associativity and unit isomorphisms between the relevant chosen pullbacks.

The unit axioms for the monoid say that the diagrams
\[
\begin{tikzcd}
C_1 \arrow[r, "{(e s,1_{C_1})}"] \arrow[dr, "1_{C_1}"']
&
C_1\times_{C_0}C_1 \arrow[d, "m"]
\\
&
C_1
\end{tikzcd}
\qquad
\begin{tikzcd}
C_1 \arrow[r, "{(1_{C_1},e t)}"] \arrow[dr, "1_{C_1}"']
&
C_1\times_{C_0}C_1 \arrow[d, "m"]
\\
&
C_1
\end{tikzcd}
\]
commute. Here the first pullback is taken over
\[
t e s=s
\qquad\text{and}\qquad
s,
\]
and the second over
\[
t
\qquad\text{and}\qquad
s e t=t.
\]
Equivalently,
\[
m(e(s(a)),a)=a,
\qquad
m(a,e(t(a)))=a.
\]

For associativity, let
\[
C_3
:=
(C_1\times_{t,C_0,s}C_1)
\times_{t\pi_2,C_0,s}
C_1.
\]
Thus a generalized element of \(C_3\) is a triple \((a,b,c)\) satisfying
\[
t(a)=s(b),
\qquad
t(b)=s(c).
\]
The monoid associativity law says that the two composites
\[
C_3\rightrightarrows C_1
\]
given by composing the first two arrows first or the last two arrows first agree:
\[
\begin{tikzcd}[column sep=large]
C_3
\arrow[r, "{m\times 1}"]
\arrow[d, "{1\times m}"']
&
C_1\times_{C_0}C_1
\arrow[d, "m"]
\\
C_1\times_{C_0}C_1
\arrow[r, "m"']
&
C_1.
\end{tikzcd}
\]
In generalized-element notation this is
\[
m(m(a,b),c)=m(a,m(b,c)).
\]
Since \(m(a,b)=b\circ a\), this is the ordinary associativity equation
\[
c\circ(b\circ a)=(c\circ b)\circ a.
\]

These are precisely the identity and associativity axioms for an internal category in \(\mathcal E\). Conversely, an internal category
\[
(C_0,C_1,s,t,e,m)
\]
determines an endospan
\[
C_0\xleftarrow{s}C_1\xrightarrow{t}C_0
\]
together with unit and multiplication maps in
\[
\operatorname{EndSpan}_{\mathcal E}(C_0).
\]
The internal category axioms are exactly the monoid axioms. Hence the two structures are equivalent.
\end{proof}

\begin{corollary}[Monads in spans are internal categories]
\label{cor:monads-in-spans-internal-categories}
Let \(\mathcal E\) be a category with pullbacks. To give a monad in \(\mathbf{Span}(\mathcal E)\) on an object \(C_0\) is equivalently to give an internal category in \(\mathcal E\) with object-of-objects \(C_0\).
\end{corollary}

\begin{proof}
By Remark~\ref{rem:monads-as-monoids-in-endohom}, monads on \(C_0\) in \(\mathbf{Span}(\mathcal E)\) are monoid objects in
\[
\operatorname{EndSpan}_{\mathcal E}(C_0).
\]
By Theorem~\ref{thm:endospan-monoids-internal-categories}, these are precisely internal categories with object-of-objects \(C_0\).
\end{proof}

\begin{definition}[Internal category]
\label{def:internal-category-as-endospan-monoid}
In this chapter, an \textbf{internal category} in \(\mathcal E\) is a monoid object in an endospan category
\[
\operatorname{EndSpan}_{\mathcal E}(C_0).
\]
Explicitly, it is a tuple
\[
\mathbb C=(C_0,C_1,s,t,e,m)
\]
where
\[
s,t:C_1\to C_0,
\qquad
e:C_0\to C_1,
\qquad
m:C_1\times_{t,C_0,s}C_1\to C_1,
\]
satisfying the source-target, identity, and associativity equations displayed in the proof of Theorem~\ref{thm:endospan-monoids-internal-categories}.
\end{definition}

\section{Mealy morphisms as monad morphisms in spans}
\label{sec:mealy-morphisms-as-monad-morphisms}

Let
\[
\mathbb A=(A_0,A_1,s_A,t_A,e_A,m_A)
\]
and
\[
\mathbb B=(B_0,B_1,s_B,t_B,e_B,m_B)
\]
be internal categories in \(\mathcal E\). We regard them as monads in \(\mathbf{Span}(\mathcal E)\):
\[
A:
A_0\xleftarrow{s_A}A_1\xrightarrow{t_A}A_0,
\qquad
B:
B_0\xleftarrow{s_B}B_1\xrightarrow{t_B}B_0.
\]

A 1-cell from \(A_0\) to \(B_0\) in \(\mathbf{Span}(\mathcal E)\) is a span
\[
A_0\xleftarrow{p}P\xrightarrow{q}B_0.
\]

There are two composite spans naturally associated with \(P\), \(A\), and \(B\).

First, the composite \(P\circ A\) is obtained by pulling back
\[
t_A:A_1\to A_0
\qquad
\text{and}
\qquad
p:P\to A_0.
\]
Thus
\[
P\circ A
=
\left(
A_0
\xleftarrow{s_A\pi_1}
A_1\times_{t_A,A_0,p}P
\xrightarrow{q\pi_2}
B_0
\right).
\]

\[
\begin{tikzcd}
A_1\times_{t_A,A_0,p}P
\arrow[r, "\pi_2"]
\arrow[d, "\pi_1"']
\arrow[dr, phantom, "\lrcorner", very near start]
&
P \arrow[d, "p"]
\\
A_1 \arrow[r, "t_A"'] & A_0.
\end{tikzcd}
\]

Second, the composite \(B\circ P\) is obtained by pulling back
\[
q:P\to B_0
\qquad
\text{and}
\qquad
s_B:B_1\to B_0.
\]
Thus
\[
B\circ P
=
\left(
A_0
\xleftarrow{p\pi_1}
P\times_{q,B_0,s_B}B_1
\xrightarrow{t_B\pi_2}
B_0
\right).
\]

\[
\begin{tikzcd}
P\times_{q,B_0,s_B}B_1
\arrow[r, "\pi_2"]
\arrow[d, "\pi_1"']
\arrow[dr, phantom, "\lrcorner", very near start]
&
B_1 \arrow[d, "s_B"]
\\
P \arrow[r, "q"'] & B_0.
\end{tikzcd}
\]

\begin{definition}[Mealy morphism]
\label{def:mealy-morphism}
Let \(\mathbb A\) and \(\mathbb B\) be internal categories, regarded as monads in \(\mathbf{Span}(\mathcal E)\). A \textbf{Mealy morphism}
\[
(P,\Phi):\mathbb A\rightsquigarrow\mathbb B
\]
is a monad morphism between these monads, with structure cell oriented as
\[
\Phi:P\circ A\Rightarrow B\circ P.
\]
Thus it consists of a span
\[
A_0\xleftarrow{p}P\xrightarrow{q}B_0
\]
and a 2-cell
\[
\Phi:
P\circ A\Rightarrow B\circ P
\]
satisfying the unit and multiplication compatibility axioms
\begin{align}
\Phi\circ(P\eta_A)&=\eta_B P,
\label{eq:mealy-unit-abstract}
\\
\Phi\circ(P\mu_A)
&=
(\mu_B P)\circ(B\Phi)\circ(\Phi A).
\label{eq:mealy-multiplication-abstract}
\end{align}
\end{definition}

\begin{remark}
\label{rem:lax-colax-convention}
Depending on variance conventions, the structure cell in Definition~\ref{def:mealy-morphism} may be called lax or colax. We shall not rely on that terminology. The displayed structure cell
\[
\Phi:P\circ A\Rightarrow B\circ P
\]
is the defining datum. Since \(\mathbf{Span}(\mathcal E)\) is a bicategory rather than a strict 2-category, the equations above suppress the associativity and unit constraints. Equivalently, the equations are read after inserting the canonical pullback isomorphisms.
\end{remark}

A 2-cell
\[
\Phi:P\circ A\Rightarrow B\circ P
\]
is a map
\[
\Phi:
A_1\times_{t_A,A_0,p}P
\longrightarrow
P\times_{q,B_0,s_B}B_1
\]
such that the diagram
\[
\begin{tikzcd}[column sep=large]
&
A_1\times_{A_0}P
\arrow[dl, "s_A\pi_1"']
\arrow[dr, "q\pi_2"]
\arrow[d, "\Phi"]
&
\\
A_0
&
P\times_{B_0}B_1
\arrow[l, "p\pi_1"]
\arrow[r, "t_B\pi_2"']
&
B_0
\end{tikzcd}
\]
commutes.

Write
\[
\Phi(a,x)=\bigl(\delta(a,x),\omega(a,x)\bigr).
\]
Then the span compatibility conditions are
\begin{align}
p(\delta(a,x))&=s_A(a),
\label{eq:mealy-left-typing}
\\
s_B(\omega(a,x))&=q(\delta(a,x)),
\label{eq:mealy-output-source}
\\
t_B(\omega(a,x))&=q(x).
\label{eq:mealy-output-target}
\end{align}

Equivalently, a Mealy morphism consists of maps
\[
\delta:
A_1\times_{t_A,A_0,p}P\to P,
\qquad
\omega:
A_1\times_{t_A,A_0,p}P\to B_1,
\]
satisfying the typing equations
\begin{align}
p\delta&=s_A\pi_1,
\label{eq:mealy-definition-pdelta}
\\
s_B\omega&=q\delta,
\label{eq:mealy-definition-source}
\\
t_B\omega&=q\pi_2.
\label{eq:mealy-definition-target}
\end{align}

\begin{proposition}[The unit law]
\label{prop:mealy-unit-law}
For a Mealy morphism \((P,\Phi):\mathbb A\rightsquigarrow\mathbb B\), the unit axiom \eqref{eq:mealy-unit-abstract} is equivalent to
\begin{align}
\delta(e_A(p x),x)&=x,
\label{eq:delta-unit}
\\
\omega(e_A(p x),x)&=e_B(qx).
\label{eq:omega-unit}
\end{align}
\end{proposition}

\begin{proof}
The unit of \(\mathbb A\) induces a map
\[
P\to A_1\times_{A_0}P,
\qquad
x\mapsto (e_A(p x),x).
\]
The unit axiom says that the composite
\[
P
\xrightarrow{(e_Ap,1_P)}
A_1\times_{A_0}P
\xrightarrow{\Phi}
P\times_{B_0}B_1
\]
equals the map
\[
P
\xrightarrow{(1_P,e_Bq)}
P\times_{B_0}B_1.
\]
Comparing the two components gives precisely
\[
\delta(e_A(p x),x)=x,
\qquad
\omega(e_A(p x),x)=e_B(qx).
\]
\end{proof}

\begin{proposition}[The multiplication law]
\label{prop:mealy-multiplication-law}
Let
\[
a:u\to v,
\qquad
b:v\to w
\]
be internally composable arrows of \(\mathbb A\), and let \(x\in P\) satisfy
\[
p(x)=w.
\]
Then the multiplication axiom \eqref{eq:mealy-multiplication-abstract} is equivalent to
\begin{align}
\delta(b\circ a,x)
&=
\delta(a,\delta(b,x)),
\label{eq:delta-composition}
\\
\omega(b\circ a,x)
&=
\omega(b,x)\circ \omega(a,\delta(b,x)).
\label{eq:omega-composition}
\end{align}
Equivalently, using \(m_A(a,b)=b\circ a\),
\begin{align}
\delta(m_A(a,b),x)
&=
\delta(a,\delta(b,x)),
\label{eq:delta-composition-m}
\\
\omega(m_A(a,b),x)
&=
m_B\bigl(\omega(a,\delta(b,x)),\omega(b,x)\bigr).
\label{eq:omega-composition-m}
\end{align}
\end{proposition}

\begin{proof}
The left-hand side of \eqref{eq:mealy-multiplication-abstract} first composes \(a\) and \(b\) in \(\mathbb A\), then applies \(\Phi\). It therefore produces
\[
\bigl(\delta(b\circ a,x),\omega(b\circ a,x)\bigr).
\]

The right-hand side processes the same data in two stages. First apply \(\Phi\) to the pair \((b,x)\). This gives
\[
\bigl(\delta(b,x),\omega(b,x)\bigr).
\]
Since
\[
p(\delta(b,x))=s_A(b)=v,
\]
the pair \((a,\delta(b,x))\) is admissible. Applying \(\Phi\) again gives
\[
\bigl(\delta(a,\delta(b,x)),\omega(a,\delta(b,x))\bigr).
\]
The two produced \(B\)-arrows are composable because
\[
t_B(\omega(a,\delta(b,x)))=q(\delta(b,x))
=
s_B(\omega(b,x)).
\]
Thus their composite in \(\mathbb B\) is
\[
\omega(b,x)\circ \omega(a,\delta(b,x)).
\]

The multiplication law asserts equality between the one-step and two-step procedures. Hence the displayed equations follow.
\end{proof}

\begin{remark}[Generalized-element notation]
\label{rem:generalized-elements}
The equations above are written in generalized-element notation. Formally, the composition equations are equalities of morphisms out of the pullback object
\[
D
:=
\bigl(A_1\times_{t_A,A_0,s_A}A_1\bigr)
\times_{t_A\pi_2,A_0,p}
P.
\]
A generalized element of \(D\) is a triple \((a,b,x)\) satisfying
\[
t_A(a)=s_A(b),
\qquad
t_A(b)=p(x).
\]
\end{remark}

\section{Operational semantics of Mealy morphisms}
\label{sec:operational-semantics-mealy}

The terminology ``Mealy morphism'' is not merely mnemonic. Classical deterministic Mealy machines are state machines in which each transition consumes an input, updates the state, and emits an output. In coalgebraic form, such a machine is often represented by a map
\[
X\to (O\times X)^I,
\]
or equivalently by a transition-output map
\[
I\times X\to X\times O.
\]
The span-internal construction replaces the untyped product \(I\times X\) by the typed pullback
\[
A_1\times_{A_0}P
\]
and replaces the output set \(O\) by the arrow object \(B_1\) of an internal category. This is the same operational idea, but internalized and typed by the object spaces of \(\mathbb A\) and \(\mathbb B\).

This viewpoint is also close in spirit to Paré's Mealy morphisms of enriched categories, which categorify Mealy machines and place them between functors and profunctors \cite{PareMealyMorphisms}. The coalgebraic treatment of Mealy machines as systems of states and observations is standard in universal coalgebra \cite{RuttenUniversalCoalgebra,BonsangueRuttenSilvaMealy}.

Let
\[
(P,\delta,\omega):\mathbb A\rightsquigarrow\mathbb B
\]
be a Mealy morphism. A state \(x\in P\) carries two interfaces:
\[
p(x)\in A_0,
\qquad
q(x)\in B_0.
\]
We write this suggestively as a state judgement
\[
x:p(x)\leadsto q(x).
\]

If
\[
a:u\to v
\]
is an arrow of \(\mathbb A\), then \(a\) is admissible at \(x\) precisely when
\[
p(x)=v.
\]
In that case the Mealy transition is
\[
x
\;\xrightarrow{\;a/\omega(a,x)\;}
\delta(a,x).
\]

The typing equations say that
\[
p(\delta(a,x))=u
\]
and that the emitted output is an arrow
\[
\omega(a,x):q(\delta(a,x))\to q(x)
\]
in \(\mathbb B\). Thus the machine processes an input arrow by moving its state backwards along the input type and emitting an output arrow between the corresponding output interfaces.

\subsection{Natural-deduction style rules}
\label{subsec:natural-deduction-mealy}

The operational calculus can be displayed as a system of typing and transition rules.

The state formation rule is
\[
\frac{x\in P}
     {x:p(x)\leadsto q(x)}.
\]

Input admissibility is expressed by
\[
\frac{
a:u\to_{\mathbb A}v
\qquad
x:v\leadsto y
}
{
a\ \text{is admissible at}\ x
}.
\]

The one-step execution rule is
\[
\frac{
a:u\to_{\mathbb A}v
\qquad
x:v\leadsto y
}
{
x
\xrightarrow{\;a/\omega(a,x)\;}
\delta(a,x)
}.
\]
The conclusion carries the additional typing information
\[
\delta(a,x):u\leadsto y',
\qquad
\omega(a,x):y'\to_{\mathbb B}y,
\]
where
\[
y'=q(\delta(a,x)).
\]

The identity rule is
\[
\frac{
x:v\leadsto y
}
{
x\xrightarrow{\;1_v/1_y\;}x
}.
\]
This is exactly the Mealy unit law
\[
\delta(e_A(p x),x)=x,
\qquad
\omega(e_A(p x),x)=e_B(qx).
\]

For sequential execution, suppose
\[
a:u\to v,
\qquad
b:v\to w,
\qquad
x:w\leadsto y.
\]
There are two ways to process the composite input
\[
b\circ a:u\to w.
\]

The one-step execution is
\[
x
\xrightarrow{\;b\circ a/\omega(b\circ a,x)\;}
\delta(b\circ a,x).
\]

The two-step execution is
\[
x
\xrightarrow{\;b/\omega(b,x)\;}
\delta(b,x)
\xrightarrow{\;a/\omega(a,\delta(b,x))\;}
\delta(a,\delta(b,x)).
\]

The multiplication law says that these two executions agree:
\[
\delta(b\circ a,x)
=
\delta(a,\delta(b,x)),
\]
and
\[
\omega(b\circ a,x)
=
\omega(b,x)\circ \omega(a,\delta(b,x)).
\]

Thus a Mealy morphism is a typed state machine whose transition and output behavior is compatible with identities and sequential composition.

\begin{remark}[Contravariant input processing]
\label{rem:contravariant-input-processing}
The input is processed target-to-source. If
\[
a:u\to v
\]
and \(p(x)=v\), then
\[
p(\delta(a,x))=u.
\]
This direction is forced by the chosen structure cell
\[
\Phi:P\circ A\Rightarrow B\circ P.
\]
Operationally, one may regard a Mealy morphism as a right action of the input category on the state span, together with an output cocycle valued in the output category.
\end{remark}

\begin{example}[The one-object case]
\label{ex:one-object-mealy}
Let \(\mathcal E=\mathbf{Set}\), and suppose \(\mathbb A\) and \(\mathbb B\) have one object. Then they are monoids, say \(M\) and \(N\). A span
\[
1\xleftarrow{}P\xrightarrow{}1
\]
is just a set \(P\). A Mealy morphism consists of functions
\[
\delta:M\times P\to P,
\qquad
\omega:M\times P\to N,
\]
satisfying
\[
\delta(1,x)=x,
\qquad
\omega(1,x)=1.
\]

To match the categorical composition convention used throughout the chapter, write monoid multiplication by
\[
mn:=m\circ n.
\]
Thus \(mn\) means first \(n\), then \(m\). With this convention the composition laws are
\[
\delta(mn,x)=\delta(n,\delta(m,x)),
\]
and
\[
\omega(mn,x)=\omega(m,x)\,\omega(n,\delta(m,x)).
\]
Equivalently, writing
\[
x\cdot m:=\delta(m,x),
\]
the state update is a right action:
\[
x\cdot(mn)=(x\cdot m)\cdot n,
\]
and \(\omega\) is the corresponding output cocycle:
\[
\omega(mn,x)=\omega(m,x)\,\omega(n,x\cdot m).
\]
\end{example}

\section{Composition of Mealy morphisms}
\label{sec:composition-of-mealy-morphisms}

Mealy morphisms compose by the ordinary composition of monad morphisms in a bicategory. Explicitly, let
\[
(P,\delta_P,\omega_P):\mathbb A\rightsquigarrow\mathbb B
\]
and
\[
(Q,\delta_Q,\omega_Q):\mathbb B\rightsquigarrow\mathbb C
\]
be Mealy morphisms, with underlying spans
\[
A_0\xleftarrow{p_P}P\xrightarrow{q_P}B_0,
\qquad
B_0\xleftarrow{p_Q}Q\xrightarrow{q_Q}C_0.
\]

Their composite has carrier span
\[
A_0
\xleftarrow{p_P\pi_P}
P\times_{q_P,B_0,p_Q}Q
\xrightarrow{q_Q\pi_Q}
C_0.
\]

Operationally, a composite state is a pair \((x,y)\) satisfying
\[
q_P(x)=p_Q(y).
\]
Given an input arrow
\[
a:u\to v
\]
of \(\mathbb A\) with
\[
p_P(x)=v,
\]
the first machine runs
\[
x
\xrightarrow{\;a/\omega_P(a,x)\;}
\delta_P(a,x).
\]
It emits a \(\mathbb B\)-arrow
\[
\omega_P(a,x):q_P(\delta_P(a,x))\to q_P(x).
\]
Since
\[
q_P(x)=p_Q(y),
\]
this emitted \(\mathbb B\)-arrow is admissible input for the second machine at \(y\). The second machine then runs
\[
y
\xrightarrow{\;\omega_P(a,x)/\omega_Q(\omega_P(a,x),y)\;}
\delta_Q(\omega_P(a,x),y).
\]

Thus the composite Mealy transition is
\[
\delta_{QP}(a,(x,y))
=
\bigl(
\delta_P(a,x),
\delta_Q(\omega_P(a,x),y)
\bigr),
\]
and the composite output is
\[
\omega_{QP}(a,(x,y))
=
\omega_Q(\omega_P(a,x),y).
\]

The identity Mealy morphism on \(\mathbb A\) has underlying identity span
\[
A_0\xleftarrow{1_{A_0}}A_0\xrightarrow{1_{A_0}}A_0.
\]
At a state \(v\in A_0\), an arrow
\[
a:u\to v
\]
is processed by
\[
v
\xrightarrow{\;a/a\;}
u.
\]
Thus the identity Mealy morphism has
\[
\delta(a,v)=u=s_A(a),
\qquad
\omega(a,v)=a.
\]

\begin{remark}
\label{rem:pipeline-composition}
Composition of Mealy morphisms is pipeline composition. The first machine consumes \(\mathbb A\)-input and emits \(\mathbb B\)-output; the second machine consumes that emitted \(\mathbb B\)-arrow and emits \(\mathbb C\)-output. The bicategorical associativity of this composition is inherited from the associativity constraints of \(\mathbf{Span}(\mathcal E)\).
\end{remark}

\section{Transformations of Mealy morphisms}
\label{sec:transformations-of-mealy-morphisms}

The bicategory \(\mathbf{Span}(\mathcal E)\) also supplies a notion of 2-cell between Mealy morphisms.

Let
\[
(P,\Phi),(Q,\Psi):\mathbb A\rightsquigarrow\mathbb B
\]
be Mealy morphisms. Write
\[
A_0\xleftarrow{p_P}P\xrightarrow{q_P}B_0,
\qquad
A_0\xleftarrow{p_Q}Q\xrightarrow{q_Q}B_0
\]
for their underlying spans.

\begin{definition}[Mealy transformation]
\label{def:mealy-transformation}
A \textbf{Mealy transformation}
\[
\theta:(P,\Phi)\Rightarrow(Q,\Psi)
\]
is a map of spans
\[
\theta:P\to Q
\]
such that
\[
p_Q\theta=p_P,
\qquad
q_Q\theta=q_P,
\]
and such that the diagram
\[
\begin{tikzcd}[column sep=large]
A_1\times_{A_0}P
\arrow[r, "\Phi"]
\arrow[d, "1_{A_1}\times\theta"']
&
P\times_{B_0}B_1
\arrow[d, "\theta\times 1_{B_1}"]
\\
A_1\times_{A_0}Q
\arrow[r, "\Psi"']
&
Q\times_{B_0}B_1
\end{tikzcd}
\]
commutes.
\end{definition}

Unpacked in components, the compatibility condition says
\begin{align}
\theta(\delta_P(a,x))
&=
\delta_Q(a,\theta x),
\label{eq:mealy-transformation-delta}
\\
\omega_P(a,x)
&=
\omega_Q(a,\theta x).
\label{eq:mealy-transformation-omega}
\end{align}

\begin{proposition}
\label{prop:mealy-transformations-compose}
For fixed internal categories \(\mathbb A\) and \(\mathbb B\), Mealy morphisms
\[
\mathbb A\rightsquigarrow\mathbb B
\]
and Mealy transformations between them form a category.
\end{proposition}

\begin{proof}
The identity transformation on \((P,\Phi)\) is the identity map
\[
1_P:P\to P.
\]
It is a map of spans and trivially commutes with the structure map \(\Phi\).

If
\[
\theta:(P,\Phi)\Rightarrow(Q,\Psi)
\]
and
\[
\kappa:(Q,\Psi)\Rightarrow(R,\Xi)
\]
are Mealy transformations, then \(\kappa\theta:P\to R\) is again a map of spans. Moreover,
\[
\Xi\circ(1_{A_1}\times \kappa\theta)
=
(\kappa\times 1_{B_1})\circ\Psi\circ(1_{A_1}\times\theta)
=
(\kappa\times 1_{B_1})\circ(\theta\times 1_{B_1})\circ\Phi
=
((\kappa\theta)\times 1_{B_1})\circ\Phi.
\]
Thus \(\kappa\theta\) is a Mealy transformation. Associativity and unit laws follow from those in \(\mathcal E\).
\end{proof}

\section{Representable spans and canonical adjunctions}
\label{sec:representable-spans-canonical-adjunctions}

A general Mealy morphism
\[
(P,\Phi):\mathbb A\rightsquigarrow\mathbb B
\]
is carried by an arbitrary span
\[
A_0\xleftarrow{p}P\xrightarrow{q}B_0.
\]
The object \(P\) is then a genuine state object. Nothing in the definition requires this span to arise from a map \(A_0\to B_0\), and nothing requires it to carry an adjoint.

The ordinary functorial case appears when the carrier span has extra structure: it is the canonical span determined by a map of object-objects.

Let
\[
F_0:A_0\to B_0
\]
be a morphism in \(\mathcal E\). It determines a span
\[
F_{0!}
:=
\left(
A_0
\xleftarrow{1_{A_0}}
A_0
\xrightarrow{F_0}
B_0
\right).
\]

It also determines the reverse span
\[
F_0^\ast
:=
\left(
B_0
\xleftarrow{F_0}
A_0
\xrightarrow{1_{A_0}}
A_0
\right).
\]

Thus
\[
F_{0!}:A_0\to B_0,
\qquad
F_0^\ast:B_0\to A_0
\]
are 1-cells in \(\mathbf{Span}(\mathcal E)\).

\[
\begin{tikzcd}[column sep=large]
A_0
&
A_0 \arrow[l, "1_{A_0}"'] \arrow[r, "F_0"]
&
B_0
\end{tikzcd}
\qquad
\begin{tikzcd}[column sep=large]
B_0
&
A_0 \arrow[l, "F_0"'] \arrow[r, "1_{A_0}"]
&
A_0.
\end{tikzcd}
\]

\begin{proposition}[Canonical adjunction associated to a map]
\label{prop:canonical-span-adjunction}
There is a canonical adjunction
\[
F_{0!}\dashv F_0^\ast
\]
in \(\mathbf{Span}(\mathcal E)\).
\end{proposition}

\begin{proof}
First compute
\[
F_0^\ast\circ F_{0!}:A_0\to A_0.
\]
Its apex is the pullback
\[
A_0\times_{B_0}A_0
\]
of \(F_0\) with itself:
\[
\begin{tikzcd}
A_0\times_{B_0}A_0
\arrow[r, "\pi_2"]
\arrow[d, "\pi_1"']
\arrow[dr, phantom, "\lrcorner", very near start]
&
A_0 \arrow[d, "F_0"]
\\
A_0 \arrow[r, "F_0"'] & B_0.
\end{tikzcd}
\]
The composite span is
\[
A_0
\xleftarrow{\pi_1}
A_0\times_{B_0}A_0
\xrightarrow{\pi_2}
A_0.
\]
The unit
\[
\eta_F:1_{A_0}\Rightarrow F_0^\ast\circ F_{0!}
\]
is induced by the diagonal map
\[
\Delta:A_0\to A_0\times_{B_0}A_0,
\]
since
\[
\pi_1\Delta=1_{A_0},
\qquad
\pi_2\Delta=1_{A_0}.
\]

Next compute
\[
F_{0!}\circ F_0^\ast:B_0\to B_0.
\]
Its apex is canonically isomorphic to \(A_0\), and under this canonical identification the composite span is
\[
B_0\xleftarrow{F_0}A_0\xrightarrow{F_0}B_0.
\]
The counit
\[
\varepsilon_F:F_{0!}\circ F_0^\ast\Rightarrow 1_{B_0}
\]
is induced by the map
\[
F_0:A_0\to B_0.
\]

The triangle identities follow from the equations defining the diagonal \(\Delta\) and the universal properties of the relevant pullbacks. Hence
\[
F_{0!}\dashv F_0^\ast.
\]
\end{proof}

\begin{definition}[Representable span]
\label{def:representable-span}
A span
\[
A_0\xleftarrow{p}P\xrightarrow{q}B_0
\]
is called \textbf{representable} if it is isomorphic, as a span, to a span of the form
\[
A_0
\xleftarrow{1_{A_0}}
A_0
\xrightarrow{F_0}
B_0
\]
for some map \(F_0:A_0\to B_0\).
\end{definition}

Equivalently, a span
\[
A_0\xleftarrow{p}P\xrightarrow{q}B_0
\]
is representable iff \(p\) is an isomorphism. In that case it is isomorphic to
\[
A_0\xleftarrow{1_{A_0}}A_0\xrightarrow{q p^{-1}}B_0.
\]

\begin{remark}
\label{rem:representable-versus-left-adjoint}
Every representable span has a canonical right adjoint after choosing a representative of the form
\[
A_0\xleftarrow{1_{A_0}}A_0\xrightarrow{F_0}B_0.
\]
In this chapter, we only use this direction. We do not need the converse characterization of left adjoint 1-cells in \(\mathbf{Span}(\mathcal E)\).
\end{remark}

\begin{definition}[Representable Mealy morphism]
\label{def:representable-mealy-morphism}
A Mealy morphism
\[
(P,\Phi):\mathbb A\rightsquigarrow\mathbb B
\]
is \textbf{representable} if its underlying carrier span is representable. Equivalently, up to isomorphism of carrier spans, it is carried by a canonical span
\[
F_{0!}
=
\left(
A_0
\xleftarrow{1_{A_0}}
A_0
\xrightarrow{F_0}
B_0
\right),
\]
which has the canonical right adjoint
\[
F_0^\ast
=
\left(
B_0
\xleftarrow{F_0}
A_0
\xrightarrow{1_{A_0}}
A_0
\right).
\]
\end{definition}

\begin{remark}
\label{rem:adjunction-as-structural-certificate}
The adjunction
\[
F_{0!}\dashv F_0^\ast
\]
is not an operation that turns an arbitrary Mealy morphism into a functor. Rather, it is additional structure available for a representable carrier. It certifies that the carrier is map-like, rather than genuinely stateful. This is the structural distinction between general Mealy morphisms and internal functors.
\end{remark}

\section{Internal functors as representable Mealy morphisms}
\label{sec:internal-functors-left-adjoint-mealy}

The previous section showed that every map
\[
F_0:A_0\to B_0
\]
determines a canonical adjunction
\[
F_{0!}\dashv F_0^\ast
\]
in \(\mathbf{Span}(\mathcal E)\). We now show that Mealy morphisms carried by such canonical spans are precisely internal functors.

Thus the functorial case is not merely the case of a simple state object. It is the case where the Mealy carrier is representable:
\[
\text{Mealy morphism}
+
\text{canonical representable carrier}
=
\text{internal functor}.
\]

Let
\[
F_0:A_0\to B_0
\]
be a morphism in \(\mathcal E\), and consider the canonical carrier
\[
F_{0!}
=
\left(
A_0
\xleftarrow{1_{A_0}}
A_0
\xrightarrow{F_0}
B_0
\right).
\]

A Mealy structure on \(F_{0!}\) is a 2-cell
\[
\Phi:F_{0!}\circ A\Rightarrow B\circ F_{0!}.
\]

We calculate this structure explicitly.

First,
\[
F_{0!}\circ A
\cong
\left(
A_0
\xleftarrow{s_A}
A_1
\xrightarrow{F_0t_A}
B_0
\right),
\]
because
\[
A_1\times_{t_A,A_0,1_{A_0}}A_0\cong A_1.
\]

Second,
\[
B\circ F_{0!}
=
\left(
A_0
\xleftarrow{\pi_1}
A_0\times_{F_0,B_0,s_B}B_1
\xrightarrow{t_B\pi_2}
B_0
\right).
\]

\[
\begin{tikzcd}
A_0\times_{F_0,B_0,s_B}B_1
\arrow[r, "\pi_2"]
\arrow[d, "\pi_1"']
\arrow[dr, phantom, "\lrcorner", very near start]
&
B_1 \arrow[d, "s_B"]
\\
A_0 \arrow[r, "F_0"'] & B_0.
\end{tikzcd}
\]

Thus \(\Phi\) is equivalently a map
\[
\Phi:A_1\to A_0\times_{B_0}B_1
\]
such that
\[
\begin{tikzcd}[column sep=large]
&
A_1
\arrow[dl, "s_A"']
\arrow[dr, "F_0t_A"]
\arrow[d, "\Phi"]
&
\\
A_0
&
A_0\times_{B_0}B_1
\arrow[l, "\pi_1"]
\arrow[r, "t_B\pi_2"']
&
B_0.
\end{tikzcd}
\]

Write
\[
\Phi(a)=\bigl(\delta(a),F_1(a)\bigr).
\]
The condition that \(\Phi(a)\) lies in the pullback says
\[
F_0(\delta(a))=s_B(F_1(a)).
\]
The condition that \(\Phi\) is a map of spans says
\[
\delta(a)=s_A(a),
\]
and
\[
t_B(F_1(a))=F_0(t_A(a)).
\]
Therefore
\[
s_BF_1=F_0s_A,
\qquad
t_BF_1=F_0t_A.
\]

The point is structural: because the carrier span is
\[
F_{0!}:A_0\to B_0,
\]
its state object is \(A_0\) itself. Hence there is no independent state update to choose. The update component of the Mealy structure is forced to be
\[
\delta=s_A.
\]
The remaining component \(F_1\) is exactly the arrow map of a functor.

\begin{definition}[Internal functor]
\label{def:internal-functor}
Let
\[
\mathbb A=(A_0,A_1,s_A,t_A,e_A,m_A)
\]
and
\[
\mathbb B=(B_0,B_1,s_B,t_B,e_B,m_B)
\]
be internal categories in \(\mathcal E\). An \textbf{internal functor}
\[
F:\mathbb A\to\mathbb B
\]
consists of maps
\[
F_0:A_0\to B_0,
\qquad
F_1:A_1\to B_1,
\]
such that
\begin{align}
s_BF_1&=F_0s_A,
\label{eq:internal-functor-source}
\\
t_BF_1&=F_0t_A,
\label{eq:internal-functor-target}
\\
F_1e_A&=e_BF_0,
\label{eq:internal-functor-identity}
\\
F_1m_A&=m_B(F_1\times F_1).
\label{eq:internal-functor-composition}
\end{align}
Here \(F_1\times F_1\) denotes the induced map on pullbacks
\[
A_1\times_{t_A,A_0,s_A}A_1
\longrightarrow
B_1\times_{t_B,B_0,s_B}B_1,
\]
which is well-defined by the source-target equations.
\end{definition}

The source-target conditions are expressed by
\[
\begin{tikzcd}
A_1 \arrow[r, "F_1"] \arrow[d, "s_A"'] & B_1 \arrow[d, "s_B"] \\
A_0 \arrow[r, "F_0"'] & B_0
\end{tikzcd}
\qquad
\begin{tikzcd}
A_1 \arrow[r, "F_1"] \arrow[d, "t_A"'] & B_1 \arrow[d, "t_B"] \\
A_0 \arrow[r, "F_0"'] & B_0.
\end{tikzcd}
\]
The identity and composition conditions are expressed by
\[
\begin{tikzcd}
A_0 \arrow[r, "F_0"] \arrow[d, "e_A"'] & B_0 \arrow[d, "e_B"] \\
A_1 \arrow[r, "F_1"'] & B_1
\end{tikzcd}
\]
and
\[
\begin{tikzcd}
A_1\times_{A_0}A_1
\arrow[r, "F_1\times F_1"]
\arrow[d, "m_A"']
&
B_1\times_{B_0}B_1
\arrow[d, "m_B"]
\\
A_1 \arrow[r, "F_1"'] & B_1.
\end{tikzcd}
\]

\begin{theorem}[Mealy morphisms carried by canonical representables are internal functors]
\label{thm:left-adjoint-mealy-internal-functor}
Let
\[
\mathbb A,\mathbb B
\]
be internal categories in \(\mathcal E\), and let
\[
F_0:A_0\to B_0
\]
be a morphism in \(\mathcal E\). To give an internal functor
\[
F:\mathbb A\to\mathbb B
\]
with object map \(F_0\) is equivalently to give a Mealy morphism
\[
(F_{0!},\Phi):\mathbb A\rightsquigarrow\mathbb B
\]
whose carrier is the canonical span
\[
F_{0!}
=
\left(
A_0
\xleftarrow{1_{A_0}}
A_0
\xrightarrow{F_0}
B_0
\right).
\]
This canonical carrier has the right adjoint
\[
F_0^\ast
=
\left(
B_0
\xleftarrow{F_0}
A_0
\xrightarrow{1_{A_0}}
A_0
\right).
\]
\end{theorem}

\begin{proof}
Suppose first that
\[
(F_{0!},\Phi):\mathbb A\rightsquigarrow\mathbb B
\]
is a Mealy morphism carried by the canonical span
\[
F_{0!}
=
\left(
A_0
\xleftarrow{1_{A_0}}
A_0
\xrightarrow{F_0}
B_0
\right).
\]
By the calculation above, \(\Phi\) is equivalently a map
\[
\Phi:A_1\to A_0\times_{B_0}B_1.
\]
Write
\[
\Phi(a)=(\delta(a),F_1(a)).
\]
Since \(\Phi\) is a map of spans, its left-leg compatibility forces
\[
\delta=s_A.
\]
The pullback condition and right-leg compatibility then give
\[
s_BF_1=F_0s_A,
\qquad
t_BF_1=F_0t_A.
\]

The Mealy unit law says
\[
\omega(e_A(x),x)=e_B(F_0x).
\]
Since in the canonical representable carrier case the output component is \(F_1\), this is exactly
\[
F_1e_A=e_BF_0.
\]

For multiplication, let
\[
a:u\to v,
\qquad
b:v\to w.
\]
The Mealy multiplication law says
\[
\omega(b\circ a,w)
=
\omega(b,w)\circ \omega(a,v).
\]
Since
\[
\omega(c,t_A(c))=F_1(c),
\]
this becomes
\[
F_1(b\circ a)=F_1(b)\circ F_1(a).
\]
Equivalently,
\[
F_1m_A=m_B(F_1\times F_1).
\]
Thus \((F_0,F_1)\) is an internal functor.

Conversely, suppose
\[
F=(F_0,F_1):\mathbb A\to\mathbb B
\]
is an internal functor. The object map \(F_0\) determines the canonical adjunction
\[
F_{0!}\dashv F_0^\ast
\]
in \(\mathbf{Span}(\mathcal E)\). Define
\[
\Phi:A_1\to A_0\times_{B_0}B_1
\]
by
\[
\Phi(a):=(s_A(a),F_1(a)).
\]
This is well-defined because
\[
F_0s_A(a)=s_BF_1(a).
\]
The left leg condition is immediate:
\[
\pi_1\Phi=s_A.
\]
The right leg condition is
\[
t_B\pi_2\Phi=t_BF_1=F_0t_A.
\]
Thus \(\Phi\) is a 2-cell
\[
F_{0!}\circ A\Rightarrow B\circ F_{0!}.
\]

The identity preservation equation
\[
F_1e_A=e_BF_0
\]
is precisely the Mealy unit law. The composition preservation equation
\[
F_1m_A=m_B(F_1\times F_1)
\]
is precisely the Mealy multiplication law. Hence \((F_{0!},\Phi)\) is a Mealy morphism carried by the canonical representable span associated to \(F_0\).

The two constructions are inverse to one another.
\end{proof}

\subsection{Operational meaning of the representable carrier}
\label{subsec:operational-meaning-left-adjoint-carrier}

The operational interpretation of internal functors is not separate from the operational interpretation of Mealy morphisms. It is the same Mealy interpretation restricted to canonical representable carriers
\[
F_{0!}\dashv F_0^\ast.
\]

For a general Mealy morphism, a state is an element
\[
x\in P
\]
with two interfaces
\[
p(x)\in A_0,
\qquad
q(x)\in B_0.
\]
The state object \(P\) may contain hidden information not determined by either interface.

For a Mealy morphism carried by
\[
F_{0!}
=
\left(
A_0
\xleftarrow{1_{A_0}}
A_0
\xrightarrow{F_0}
B_0
\right),
\]
the state object is \(A_0\) itself. A state is just an object \(v\in A_0\), and its output interface is forced to be \(F_0v\). The state judgement becomes
\[
v:v\leadsto F_0v.
\]

The right adjoint
\[
F_0^\ast
=
\left(
B_0
\xleftarrow{F_0}
A_0
\xrightarrow{1_{A_0}}
A_0
\right)
\]
records the same object map in the reverse span direction. Thus the adjunction
\[
F_{0!}\dashv F_0^\ast
\]
is the bicategorical certificate that the carrier is not an arbitrary state span but a map-like, representable carrier. Operationally, this is what removes hidden state.

An input arrow
\[
a:u\to v
\]
is admissible at the state \(v\), and the updated state is forced to be \(u\). The transition rule is
\[
\frac{
a:u\to_{\mathbb A}v
}
{
v
\xrightarrow{\;a/F_1(a)\;}
u
}.
\]
The output arrow has type
\[
F_1(a):F_0u\to_{\mathbb B}F_0v.
\]

The identity rule becomes
\[
\frac{
v\in A_0
}
{
v
\xrightarrow{\;1_v/1_{F_0v}\;}
v
},
\]
which is exactly
\[
F_1(1_v)=1_{F_0v}.
\]

The sequential execution rule says that, for
\[
a:u\to v,
\qquad
b:v\to w,
\]
the one-step execution
\[
w
\xrightarrow{\;b\circ a/F_1(b\circ a)\;}
u
\]
agrees with the two-step execution
\[
w
\xrightarrow{\;b/F_1(b)\;}
v
\xrightarrow{\;a/F_1(a)\;}
u.
\]
Thus
\[
F_1(b\circ a)=F_1(b)\circ F_1(a).
\]

So an internal functor is a Mealy morphism whose carrier is the canonical representable span associated to an object map. The adjunction is what marks the stateless regime of the Mealy calculus: the state is the current input object, the update is forced, and the emitted arrow is functorial transport.

\section{Internal categories and internal functors inside the Mealy calculus}
\label{sec:internal-functors-inside-mealy}

The constructions above give the intended hierarchy.

First, internal categories are monoids in endospan categories:
\[
\mathbb A
=
(A_0,A_1,s_A,t_A,e_A,m_A)
\]
is a monoid object in
\[
\operatorname{EndSpan}_{\mathcal E}(A_0).
\]
Thus an internal category is already a typed transition system with compositional traces.

Second, Mealy morphisms are monad morphisms between these monoids in the ambient bicategory of spans:
\[
(P,\Phi):\mathbb A\rightsquigarrow\mathbb B.
\]
They are stateful implementations between internal categories. Their carrier
\[
A_0\xleftarrow{p}P\xrightarrow{q}B_0
\]
may be an arbitrary state span.

Third, internal functors appear when this carrier is representable. Namely, an object map
\[
F_0:A_0\to B_0
\]
determines
\[
F_{0!}\dashv F_0^\ast.
\]
A Mealy morphism carried by \(F_{0!}\) is exactly an internal functor.

Hence the ordinary category of internal categories and internal functors is not external to the Mealy calculus. It is the subcalculus of Mealy morphisms whose carriers are representable, with canonical representatives of the form
\[
F_{0!}
=
\left(
A_0
\xleftarrow{1_{A_0}}
A_0
\xrightarrow{F_0}
B_0
\right).
\]

\begin{theorem}[Internal functors as the representable Mealy subcalculus]
\label{thm:internal-functors-left-adjoint-subcalculus}
The ordinary category of internal categories and internal functors identifies, up to the canonical span isomorphisms inherent in \(\mathbf{Span}(\mathcal E)\), with the subcalculus of Mealy morphisms carried by canonical representable spans
\[
F_{0!}
=
\left(
A_0
\xleftarrow{1_{A_0}}
A_0
\xrightarrow{F_0}
B_0
\right).
\]
Equivalently, the essential image consists of the Mealy morphisms whose underlying carrier span is representable, i.e. whose left leg is an isomorphism.

Under this identification, an internal functor
\[
F:\mathbb A\to\mathbb B
\]
is interpreted as the stateless Mealy morphism carried by the canonical adjunction
\[
F_{0!}\dashv F_0^\ast.
\]
\end{theorem}

\begin{proof}
By Proposition~\ref{prop:canonical-span-adjunction}, every span of the form
\[
F_{0!}
=
\left(
A_0
\xleftarrow{1_{A_0}}
A_0
\xrightarrow{F_0}
B_0
\right)
\]
has a canonical right adjoint
\[
F_0^\ast
=
\left(
B_0
\xleftarrow{F_0}
A_0
\xrightarrow{1_{A_0}}
A_0
\right).
\]
By Theorem~\ref{thm:left-adjoint-mealy-internal-functor}, a Mealy structure carried by \(F_{0!}\) is exactly an internal functor with object map \(F_0\).

Composition is preserved up to the canonical span isomorphisms because composition of canonical representable spans is canonically isomorphic to a canonical representable span:
\[
G_{0!}\circ F_{0!}\cong (G_0F_0)_!.
\]
The corresponding right adjoint is
\[
F_0^\ast\circ G_0^\ast,
\]
as expected for composite adjunctions. Operationally, pipeline composition of two stateless Mealy machines is again stateless, and the resulting output map is the composite internal functor.

If one works with chosen normal forms for representable spans, this gives a literal subcategory. In invariant bicategorical terms, its essential image is the class of representable Mealy morphisms.
\end{proof}

\begin{remark}
\label{rem:automata-theoretic-internal-categories}
In this sense, internal category theory appears as a special automata theory. An internal category is a typed transition system whose transitions compose. A Mealy morphism is a stateful implementation between such transition systems. An internal functor is the special case in which the implementation is carried by a canonical representable span; equivalently, the implementation has no hidden state, the state is the current input object, and the emitted output arrow is functorial transport.
\end{remark}

\section{The main consolidation}
\label{sec:main-consolidation}

We now collect the constructions and identifications of the chapter.

\begin{theorem}[Span construction, internal categories, and Mealy morphisms]
\label{thm:main-consolidation}
Let \(\mathcal E\) be a category with pullbacks.

\begin{enumerate}
\item The pullback construction of spans determines a bicategory
\[
\mathbf{Span}(\mathcal E).
\]

\item For each object \(C_0\in\mathcal E\), the hom-category
\[
\mathbf{Span}(\mathcal E)(C_0,C_0)
\]
is a monoidal category
\[
\operatorname{EndSpan}_{\mathcal E}(C_0).
\]

\item Monads on \(C_0\) in \(\mathbf{Span}(\mathcal E)\), equivalently monoid objects in
\[
\operatorname{EndSpan}_{\mathcal E}(C_0),
\]
are precisely internal categories in \(\mathcal E\) with object-of-objects \(C_0\).

\item Mealy morphisms
\[
(P,\Phi):\mathbb A\rightsquigarrow\mathbb B
\]
are precisely monad morphisms in \(\mathbf{Span}(\mathcal E)\) with structure cell
\[
\Phi:P\circ A\Rightarrow B\circ P.
\]

\item In components, a Mealy morphism is a typed internal Mealy machine
\[
A_1\times_{A_0}P
\longrightarrow
P\times_{B_0}B_1,
\]
which consumes arrows of \(\mathbb A\), updates a state in \(P\), and emits arrows of \(\mathbb B\), compatibly with identities and sequential composition.

\item Every object map
\[
F_0:A_0\to B_0
\]
determines a canonical adjunction
\[
F_{0!}\dashv F_0^\ast
\]
in \(\mathbf{Span}(\mathcal E)\).

\item A Mealy morphism carried by such a canonical representable span
\[
F_{0!}
=
\left(
A_0
\xleftarrow{1_{A_0}}
A_0
\xrightarrow{F_0}
B_0
\right)
\]
is precisely an internal functor
\[
F:\mathbb A\to\mathbb B
\]
with object map \(F_0\).

\item Consequently, internal categories and internal functors form the representable, stateless subcalculus of the broader Mealy calculus. In invariant terms, the essential image consists of the Mealy morphisms whose carrier spans have invertible left leg.
\end{enumerate}
\end{theorem}

\begin{proof}
Part (1) is Proposition~\ref{prop:span-is-bicategory}. Part (2) follows from the bicategorical hom-category structure. Part (3) is Theorem~\ref{thm:endospan-monoids-internal-categories} together with Corollary~\ref{cor:monads-in-spans-internal-categories}. Part (4) is Definition~\ref{def:mealy-morphism}. Part (5) is the component calculation and the operational interpretation of Sections~\ref{sec:mealy-morphisms-as-monad-morphisms} and~\ref{sec:operational-semantics-mealy}. Part (6) is Proposition~\ref{prop:canonical-span-adjunction}. Part (7) is Theorem~\ref{thm:left-adjoint-mealy-internal-functor}. Part (8) is Theorem~\ref{thm:internal-functors-left-adjoint-subcalculus}.
\end{proof}

\begin{remark}
\label{rem:stateful-versus-functorial}
The bicategory of spans therefore contains two levels of morphism between internal categories. The general Mealy morphism is stateful: it processes arrows of \(\mathbb A\), updates a state lying over the object spaces of \(\mathbb A\) and \(\mathbb B\), and emits arrows of \(\mathbb B\).

Internal functors arise when the Mealy carrier is not an arbitrary span but a canonical representable span associated to an object map:
\[
F_{0!}\dashv F_0^\ast.
\]
This additional representable structure is what marks the stateless functorial regime inside the broader Mealy calculus. In that regime, the state object has collapsed to \(A_0\), the update is forced by source, and the only remaining behavior is the functorial transport of arrows.
\end{remark}